\begin{table}[htbp]\centering\small
\caption{하방비중과 삼차 변수의 양방향 예측 비교}\label{tab:downside}
\begin{tabular}{lrrrrrr}\toprule
종목 & $\beta$ & $h$ & MD$|$M1 & MDK$|$MD & MDK$|$M2 & MD$|$M2\\\midrule
\multicolumn{7}{l}{\textit{상장 전: 목표 구간 전체가 상장 전}}\\
삼성전자 & $+2$ & 5 & 2.77 & -0.28 & 0.32 & 0.60\\
삼성전자 & $-2$ & 5 & 2.76 & -0.30 & 0.32 & 0.62\\
삼성전자 & $+2$ & 20 & 6.00 & -2.04 & -0.52 & 1.49\\
삼성전자 & $-2$ & 20 & 5.97 & -1.69 & 0.04 & 1.70\\
SK하이닉스 & $+2$ & 5 & 0.47 & -1.09 & -1.08 & 0.01\\
SK하이닉스 & $-2$ & 5 & 0.45 & -1.12 & -1.07 & 0.05\\
SK하이닉스 & $+2$ & 20 & 1.87 & -1.17 & -0.57 & 0.59\\
SK하이닉스 & $-2$ & 20 & 1.54 & -1.18 & -0.58 & 0.59\\
\midrule
\multicolumn{7}{l}{\textit{상장 후: 예측 시점이 상장일 이후}}\\
삼성전자 & $+2$ & 5 & -12.49 & -0.08 & -8.87 & -8.78\\
삼성전자 & $-2$ & 5 & -14.13 & -0.60 & -10.09 & -9.43\\
삼성전자 & $+2$ & 20 & -18.64 & -4.85 & -20.41 & -14.84\\
삼성전자 & $-2$ & 20 & -22.63 & -6.22 & -23.36 & -16.14\\
SK하이닉스 & $+2$ & 5 & -4.24 & -2.17 & -7.12 & -4.84\\
SK하이닉스 & $-2$ & 5 & -5.44 & -2.99 & -8.66 & -5.50\\
SK하이닉스 & $+2$ & 20 & -4.25 & -5.53 & -12.15 & -6.27\\
SK하이닉스 & $-2$ & 20 & -7.86 & -6.93 & -15.49 & -8.00\\
\bottomrule\end{tabular}
\smallskip
\begin{minipage}{0.97\textwidth}\footnotesize 개선율 단위는 \%. A$|$B는 B 대비 A의 평균제곱예측오차 개선율이며 음수는 악화다. MDK$|$MD와 MDK$|$M2는 추가 정보의 방향을 바꾼 비교다. MD$|$M2는 비중첩 모형의 직접 비교이며, 이 표에 새로운 유의성 검정은 붙이지 않았다. M2$|$M1은 표 \ref{tab:level} 및 재현 CSV에 있다. 상장 전 표본 수는 $h=5$에서 1,904개, $h=20$에서 1,874개이며 상장 후는 각각 111개와 96개다.\end{minipage}
\end{table}
